\documentclass[11pt,a4paper]{article}
\usepackage[margin=25mm]{geometry}
\usepackage{amsmath,amssymb}
\usepackage{fontspec}
\setmainfont{SimSun}
\setsansfont{Microsoft YaHei}
\setmonofont{Consolas}
\XeTeXlinebreaklocale "zh"
\XeTeXlinebreakskip = 0pt plus 1pt
\usepackage{graphicx}
\usepackage{float}
\floatplacement{figure}{H}
\renewcommand{\figurename}{图}
\renewcommand{\tablename}{表}
\setcounter{secnumdepth}{-1}
\usepackage{longtable,booktabs,array}
\usepackage{hyperref}
\hypersetup{colorlinks=true,linkcolor=black,urlcolor=blue}
\setlength{\parindent}{2em}
\setlength{\parskip}{4pt}
\linespread{1.2}
\widowpenalty=10000
\clubpenalty=10000
\setkeys{Gin}{width=\linewidth,keepaspectratio}
\providecommand{\tightlist}{\setlength{\itemsep}{0pt}\setlength{\parskip}{0pt}}
\providecommand{\passthrough}[1]{#1}
\title{补充材料：选择后误差监测的开发实验与完整性审计}
\author{}
\date{2026年9月17日}
\begin{document}
\maketitle
\begin{center}与论文阶段稿配套；独立验证结果尚未并入\end{center}
\hypertarget{s1-ux4fe1ux606fux53efux89c1ux6027ux4e0eux6570ux636eux9694ux79bb}{%
\section{S1
信息可见性与数据隔离}\label{s1-ux4fe1ux606fux53efux89c1ux6027ux4e0eux6570ux636eux9694ux79bb}}

优化器只读取截至上一轮的已付费标签。每次评价前保存并刷新所选点预测、上界、已知最优、目标尺度、模式和随机状态；标签返回后才更新历史与监测。图1给出优化与诊断两条路径。诊断不向优化器提供反馈，候选池标签不作为免费训练数据。

\begin{figure}
\centering
\includegraphics{../results/supplement_v1/information_flow.png}
\caption{算法与诊断的信息路径。PRE日志先于真实标签保存，离线诊断读取冻结状态。虚构廉价标签或事后重新计算训练残差均不属于本实验。}
\end{figure}

\hypertarget{s1.1-ux4e00ux8f6eux6267ux884cux7684ux660eux786eux987aux5e8f}{%
\subsection{S1.1
一轮执行的明确顺序}\label{s1.1-ux4e00ux8f6eux6267ux884cux7684ux660eux786eux987aux5e8f}}

\begin{enumerate}
\def\labelenumi{\arabic{enumi}.}
\tightlist
\item
  读取历史，按预定规则截断训练子集；拟合RBF并保存训练索引与标签尺度。
\item
  生成256候选，预测均值；从原预测残差计算G/T/S/ST上界与支持回退。
\item
  读取已完成监测窗口与停滞次数，计算探索概率；仅由当前可见信息选择一个点。
\item
  保存候选检查点及所选点PRE记录，刷新文件；随后调用目标函数一次。
\item
  写LABEL记录，以原预测及原尺度构成残差；更新监测、停滞、历史。
\item
  达到预算即停止；已知最优值仅在结束后的统计中读取。
\end{enumerate}

启动阶段前D+10次训练、后D+10次在线预测评价，全部付费；不将训练残差充作监测残差。内部方法共享初始化，后续独立轨迹即使整数种子相同也不拥有同一候选池。

\hypertarget{s2-ux51bbux7ed3ux5b9eux73b0ux53c2ux6570}{%
\section{S2
冻结实现参数}\label{s2-ux51bbux7ed3ux5b9eux73b0ux53c2ux6570}}

\par\noindent\begin{minipage}{\linewidth}

\begin{longtable}[]{@{}ll@{}}
\caption{开发v1参数（实现选择，不作为创新）}\tabularnewline
\toprule
参数 & 数值 \\
\midrule
\endfirsthead
\toprule
参数 & 数值 \\
\midrule
\endhead
kernel\_scale & 0.25 \\
ridge & 1e-06 \\
max\_train & 256.0 \\
elite\_train & 128.0 \\
parents & 64.0 \\
candidates & 256.0 \\
F & 0.5 \\
CR & 0.9 \\
pbest & 0.2 \\
residual\_window & 100.0 \\
half\_life & 25.0 \\
spatial\_k & 20.0 \\
min\_support & 20.0 \\
monitor\_window & 30.0 \\
level & 0.9 \\
duplicate\_tol & 1e-10 \\
tau\_factor & 1e-06 \\
\bottomrule
\end{longtable}

\end{minipage}\par

标签尺度使用当前所选训练子集的总体标准差，为零时置1。距离为归一化欧氏距离。分位数逆经验CDF不插值，负上侧分位数截为零。反馈窗口包括启动在线预测记录。并列按原索引，重复容差与越界反射规则固定；额外Sobol补齐不获取目标标签。

开发实验B=10D，D=10/20时预算100/200，完整初始化40/60。每轨迹检查点分别41/50/75/100和61/100/150/200；首个检查点经过初始化修正，不能称25\%预算。诊断每状态256额外标签，所有校准器和策略复用同一组隐藏标签，不重复收集付费优化数据。

\hypertarget{s3-ux58f0ux660eux652fux6301ux4e0eux9519ux8befux6570ux91cf}{%
\section{S3
声明支持与错误数量}\label{s3-ux58f0ux660eux652fux6301ux4e0eux9519ux8befux6570ux91cf}}

图2公开各方法声明总数和真假改善数量。声明不是概率保证，上界失效目标10\%不意味着声明子集90\%成功。零声明时FDP为NA，不能写成零错误率证明可靠。

\begin{figure}
\centering
\includegraphics{../results/supplement_v1/declarations.png}
\caption{开发闭环中的改善声明计数。每维每方法72运行；柱高是声明数，颜色表示实际改善与错误声明。PEM-T在20维没有声明，图中零柱不是可靠性优势。}
\end{figure}

只有原型的受监测声明使用当前原预测上界，数值阈值为当前尺度乘1e-6。所有假声明都经日志核查满足上界失效事件包含关系；该包含关系不控制FDP分母，也不保证每轮成功概率。

\hypertarget{s4-ux53efux8ffdux6eafux7684ux5355ux8f68ux8ff9ux4f8bux5b50}{%
\section{S4
可追溯的单轨迹例子}\label{s4-ux53efux8ffdux6eafux7684ux5355ux8f68ux8ff9ux4f8bux5b50}}

事后展示f20、D=10、实例1、重复0。示例选择发生在开发结果汇总之后，未按单次最大差距挑选；它是描述性例子，不能估计典型性或替代正文8函数总体统计。两方法初始化与初始误差相同；最终PEM-T误差稍高，不能据一条轨迹判断机制普遍失效。

\par\noindent\begin{minipage}{\linewidth}

\begin{longtable}[]{@{}llllll@{}}
\caption{固定展示轨迹，后初始化阶段60次评价}\tabularnewline
\toprule
方法 & 终点误差 & 探索次数 & 上界失效次数 & 声明数 & 错误数 \\
\midrule
\endfirsthead
\toprule
方法 & 终点误差 & 探索次数 & 上界失效次数 & 声明数 & 错误数 \\
\midrule
\endhead
fixed01 & 3.201643 & 6 & 2 & 0 & 0 \\
PEM\_T & 3.514912 & 19 & 5 & 0 & 0 \\
\bottomrule
\end{longtable}

\end{minipage}\par

\begin{figure}
\centering
\includegraphics{../results/supplement_v1/canonical_case.png}
\caption{f20、10维、实例1、重复0的真实轨迹。上：当前最优误差；中：评价前探索概率；下：实际空间探索评价。虚线为初始化结束；后续差异来自独立闭环，不能解释为同池选点差异。}
\end{figure}

原始轨迹目录为\texttt{results/development\_v1/PEM\_T\_f20\_d10\_i1\_r0\_b100}与对应fixed01目录。所展示数据保存于\texttt{results/supplement\_v1/case.json}，图形由标签日志和PRE模式生成，没有平滑或删除评价。

\hypertarget{s5-ux5b8cux6574ux6027ux4e0eux9636ux6bb5ux8fb9ux754c}{%
\section{S5
完整性与阶段边界}\label{s5-ux5b8cux6574ux6027ux4e0eux9636ux6bb5ux8fb9ux754c}}

开发1,872轨迹、280,800算法评价；共同历史432轨迹、64,800算法评价；其中288轨迹为保存完整候选池而重复执行，不增加统计样本。诊断1,728状态、442,368标签，与算法隔离。所有完成标记、预算计数、初始化配对及标签顺序已核查。两组144/144轨迹等价性支持当前实现的代数分析。

统计单位为函数，重复用于实例内中位数；正文采用8函数配对比较、20项Holm校正及配对分层重采样。开发集用于选定固定概率，不能将该比较写成独立确认。软件锁、硬件记录和原型指纹分别保存在docs下，异常环境安装记录保留。

本材料不包含正在运行的新机制验证标签。新阶段先冻结协议与统计脚本，采用全24函数、10/20/40维、实例11--13；结果应另行审计、单列，不与开发数据合并。原计划实例6--10仍未读取。外部竞争方法的完整性能复现、CEC和应用未完成。

作者代码预算审查的TuRBO控制流反例另存于\texttt{results/author\_code\_audit/}，不与BBOB优化结果混合。完整性检查不能替代算法有效性；开发门槛失败如实报告。
\end{document}
